\documentclass{amsart}
% For dealing with references we use the comment environment
\usepackage{verbatim}
\newenvironment{reference}{\comment}{\endcomment}
%\newenvironment{reference}{}{}
\newenvironment{slogan}{\comment}{\endcomment}
\newenvironment{history}{\comment}{\endcomment}

% For commutative diagrams we use Xy-pic
\usepackage[all]{xy}

% We use 2cell for 2-commutative diagrams.
\xyoption{2cell}
\UseAllTwocells

% We use multicol for the list of chapters between chapters
\usepackage{multicol}

% This is generally recommended for better output
\usepackage{lmodern}
\usepackage[T1]{fontenc}

% For cross-file-references

% Package for hypertext links:
\usepackage{hyperref}

% For any local file, say "hello.tex" you want to link to please

% Theorem environments.
%
\theoremstyle{plain}
\newtheorem{theorem}[subsection]{Theorem}
\newtheorem{proposition}[subsection]{Proposition}
\newtheorem{lemma}[subsection]{Lemma}

\theoremstyle{definition}
\newtheorem{definition}[subsection]{Definition}
\newtheorem{example}[subsection]{Example}
\newtheorem{exercise}[subsection]{Exercise}
\newtheorem{situation}[subsection]{Situation}

\theoremstyle{remark}
\newtheorem{remark}[subsection]{Remark}
\newtheorem{remarks}[subsection]{Remarks}

\numberwithin{equation}{subsection}

% Macros
%
\def\lim{\mathop{\mathrm{lim}}\nolimits}
\def\colim{\mathop{\mathrm{colim}}\nolimits}
\def\Spec{\mathop{\mathrm{Spec}}}
\def\Hom{\mathop{\mathrm{Hom}}\nolimits}
\def\Ext{\mathop{\mathrm{Ext}}\nolimits}
\def\SheafHom{\mathop{\mathcal{H}\!\mathit{om}}\nolimits}
\def\SheafExt{\mathop{\mathcal{E}\!\mathit{xt}}\nolimits}
\def\Sch{\mathit{Sch}}
\def\Mor{\mathop{\mathrm{Mor}}\nolimits}
\def\Ob{\mathop{\mathrm{Ob}}\nolimits}
\def\Sh{\mathop{\mathit{Sh}}\nolimits}
\def\NL{\mathop{N\!L}\nolimits}
\def\CH{\mathop{\mathrm{CH}}\nolimits}
\def\proetale{{pro\text{-}\acute{e}tale}}
\def\etale{{\acute{e}tale}}
\def\QCoh{\mathit{QCoh}}
\def\Ker{\mathop{\mathrm{Ker}}}
\def\Im{\mathop{\mathrm{Im}}}
\def\Coker{\mathop{\mathrm{Coker}}}
\def\Coim{\mathop{\mathrm{Coim}}}

% Boxtimes
%
\DeclareMathSymbol{\boxtimes}{\mathbin}{AMSa}{"02}

%
% Macros for moduli stacks/spaces
%
\def\QCohstack{\mathcal{QC}\!\mathit{oh}}
\def\Cohstack{\mathcal{C}\!\mathit{oh}}
\def\Spacesstack{\mathcal{S}\!\mathit{paces}}
\def\Quotfunctor{\mathrm{Quot}}
\def\Hilbfunctor{\mathrm{Hilb}}
\def\Curvesstack{\mathcal{C}\!\mathit{urves}}
\def\Polarizedstack{\mathcal{P}\!\mathit{olarized}}
\def\Complexesstack{\mathcal{C}\!\mathit{omplexes}}
% \Pic is the operator that assigns to X its picard group, usage \Pic(X)
% \Picardstack_{X/B} denotes the Picard stack of X over B
% \Picardfunctor_{X/B} denotes the Picard functor of X over B
\def\Pic{\mathop{\mathrm{Pic}}\nolimits}
\def\Picardstack{\mathcal{P}\!\mathit{ic}}
\def\Picardfunctor{\mathrm{Pic}}
\def\Deformationcategory{\mathcal{D}\!\mathit{ef}}

\setlength{\emergencystretch}{3em}
\begin{document}
\title[Stacks Project, Tag 0AJE]{724 --- Quasi-compact quasi-separated formal spaces are filtered colimits of thickenings}
\maketitle
\noindent Copyright (C) 2005--2025 Johan de Jong. Stacks Project authors. Source: \href{https://stacks.math.columbia.edu/tag/0AJE}{Tag 0AJE}.

\noindent Editorial difficulty note (not author text). Difficulty H2: The author constructs invariant scheme-theoretic images inside an etale presentation, forms restricted etale equivalence relations, and descends each image to an algebraic-space thickening. Managing invariant images and their colimit is an advanced formal-geometric argument..

\noindent Editorial provenance note (not author text). History: excerpt from revision \texttt{a0444\allowbreak 6e57e\allowbreak c1fbc\allowbreak 252a8\allowbreak 71afc\allowbreak ec775\allowbreak 2fb28\allowbreak 07b14}, formal-spaces.tex, lines 3648--3728. Reference presentation was adapted; original-author context and core support blocks were retained or added. The mathematical wording is unchanged. Licensed under GNU FDL 1.2 or later; no invariant sections or cover texts. See ../licenses/COPYING. Context blocks reproduce author definitions and supporting results; they are not additional counted problems.

\begin{definition}
\label{definition-formal-algebraic-space}
Let $S$ be a scheme. We say a sheaf $X$ on $(\Sch/S)_{fppf}$ is a
{\it formal algebraic space} if there exist a family of maps
$\{X_i \to X\}_{i \in I}$ of sheaves such that
\begin{enumerate}
\item $X_i$ is an affine formal algebraic space,
\item $X_i \to X$ is representable by algebraic spaces and \'etale,
\item $\coprod X_i \to X$ is surjective as a map of sheaves
\end{enumerate}
and $X$ satisfies a set theoretic condition
(see Remark \href{https://stacks.math.columbia.edu/tag/0AIS}{remark set theoretic (Tag 0AIS)}). A
{\it morphism of formal algebraic spaces}
over $S$ is a map of sheaves.
\end{definition}

\begin{definition}
\label{definition-affine-formal-algebraic-space}
Let $S$ be a scheme. We say a sheaf $X$ on $(\Sch/S)_{fppf}$ is an
{\it affine formal algebraic space} if there exist
\begin{enumerate}
\item a directed set $\Lambda$,
\item a system $(X_\lambda, f_{\lambda \mu})$ over $\Lambda$
in $(\Sch/S)_{fppf}$ where
\begin{enumerate}
\item each $X_\lambda$ is affine,
\item each $f_{\lambda \mu} : X_\lambda \to X_\mu$ is a thickening,
\end{enumerate}
\end{enumerate}
such that
$$
X \cong \colim_{\lambda \in \Lambda} X_\lambda
$$
as fppf sheaves and $X$ satisfies a set theoretic condition
(see Remark \href{https://stacks.math.columbia.edu/tag/0AIS}{remark set theoretic (Tag 0AIS)}). A
{\it morphism of affine formal algebraic spaces}
over $S$ is a map of sheaves.
\end{definition}

\begin{lemma}
\label{lemma-structure-quasi-compact-quasi-separated}
\begin{reference}
\href{https://stacks.math.columbia.edu/bibliography}{Bibliography: Yasuda (Proposition 3.32)}
\end{reference}
Let $S$ be a scheme. Let $X$ be a quasi-compact and quasi-separated
formal algebraic space over $S$. Then $X = \colim X_\lambda$
for a system of algebraic spaces $(X_\lambda, f_{\lambda \mu})$
over a directed set $\Lambda$ where each
$f_{\lambda \mu} : X_\lambda \to X_\mu$ is a thickening.
\end{lemma}

\begin{proof}
By Lemma \href{https://stacks.math.columbia.edu/tag/0AJ9}{lemma characterize quasi compact (Tag 0AJ9)} we may choose an
affine formal algebraic space $Y$ and a representable surjective
\'etale morphism $Y \to X$. Write $Y = \colim Y_\lambda$ as in
Definition \ref{definition-affine-formal-algebraic-space}.

\medskip\noindent
Pick $\lambda \in \Lambda$. Then $Y_\lambda \times_X Y$ is a scheme by
Lemma \href{https://stacks.math.columbia.edu/tag/0AIG}{lemma presentation representable (Tag 0AIG)}. The reduction
(Lemma \href{https://stacks.math.columbia.edu/tag/0AIN}{lemma reduction formal algebraic space (Tag 0AIN)})
of $Y_\lambda \times_X Y$ is equal to the reduction of
$Y_{red} \times_{X_{red}} Y_{red}$ which is quasi-compact as $X$
is quasi-separated and $Y_{red}$ is affine.
Therefore $Y_\lambda \times_X Y$ is a quasi-compact scheme.
Hence there exists a $\mu \geq \lambda$ such that
$\text{pr}_2 : Y_\lambda \times_X Y \to Y$ factors
through $Y_\mu$, see Lemma \href{https://stacks.math.columbia.edu/tag/0AIA}{lemma factor through thickening (Tag 0AIA)}.
Let $Z_\lambda$ be the scheme theoretic image of the morphism
$\text{pr}_2 : Y_\lambda \times_X Y \to Y_\mu$.
This is independent of the choice of $\mu$ and we can and
will think of $Z_\lambda \subset Y$ as the scheme theoretic
image of the morphism $\text{pr}_2 : Y_\lambda \times_X Y \to Y$.
Observe that $Z_\lambda$ is also equal to the scheme theoretic image
of the morphism $\text{pr}_1 : Y \times_X Y_\lambda \to Y$ since
this is isomorphic to the morphism used to define $Z_\lambda$.
We claim that $Z_\lambda \times_X Y = Y \times_X Z_\lambda$ as subfunctors
of $Y \times_X Y$. Namely, since $Y \to X$ is \'etale we see that
$Z_\lambda \times_X Y$ is the scheme theoretic image of the morphism
$$
\text{pr}_{13} = \text{pr}_1 \times \text{id}_Y :
Y \times_X Y_\lambda \times_X Y \longrightarrow Y \times_X Y
$$
by Morphisms of Spaces, Lemma
\href{https://stacks.math.columbia.edu/tag/082Z}{lemma quasi compact scheme theoretic image (Tag 082Z)}.
By the same token, $Y \times_X Z_\lambda$ is the scheme theoretic image
of the morphism
$$
\text{pr}_{13} = \text{id}_Y \times \text{pr}_2 : 
Y \times_X Y_\lambda \times_X Y \longrightarrow Y \times_X Y
$$
The claim follows. Then
$R_\lambda = Z_\lambda \times_X Y = Y \times_X Z_\lambda$
together with the morphism $R_\lambda \to Z_\lambda \times_S Z_\lambda$
defines an \'etale equivalence relation. In this way we obtain an algebraic
space $X_\lambda = Z_\lambda/R_\lambda$. By construction the diagram
$$
\xymatrix{
Z_\lambda \ar[r] \ar[d] & Y \ar[d] \\
X_\lambda \ar[r] & X
}
$$
is cartesian (because $X$ is the coequalizer of the two projections
$R = Y \times_X Y \to Y$, because $Z_\lambda \subset Y$ is $R$-invariant,
and because $R_\lambda$ is the restriction of $R$ to $Z_\lambda$).
Hence $X_\lambda \to X$ is representable and a closed immersion, see
Spaces, Lemma
\href{https://stacks.math.columbia.edu/tag/03I2}{lemma morphism sheaves with P effective descent etale (Tag 03I2)}.
On the other hand, since $Y_\lambda \subset Z_\lambda$ we see that
$(X_\lambda)_{red} = X_{red}$, in other words, $X_\lambda \to X$
is a thickening. Finally, we claim that
$$
X = \colim X_\lambda
$$
We have $Y \times_X X_\lambda = Z_\lambda \supset Y_\lambda$. Every
morphism $T \to X$ where $T$ is a scheme over $S$ lifts \'etale locally
to a morphism into $Y$ which lifts \'etale locally into a morphism
into some $Y_\lambda$. Hence $T \to X$ lifts \'etale locally on
$T$ to a morphism into $X_\lambda$. This finishes the proof.
\end{proof}
\end{document}
